\input{../tex-inputs/latex-leseansicht-vorspann}

\section[Arthur Schnitzler an Rainer Maria Rilke, 7. 12. 1901]{L04344 Arthur Schnitzler an Rainer Maria Rilke, 7. 12. 1901}
\nopagebreak\mylabel{L04344v}
\rehead{ }\normalsize\beginnumbering\briefempfaengerindex{Rilke, Rainer Maria@\textsc{Rilke, Rainer Maria}!zzzSchnitzler, Arthur@\emph{von Arthur Schnitzler}!1901-12-071@{7. 12. 1901}|(be}
\toendnotes[C]{\smallbreak\pagebreak[2]}
\correspDesc{Versand  durch Arthur Schnitzler am 7. 12. 1901 in Wien
\newline{}Erhalt  durch Rainer Maria Rilke im Zeitraum [8. 12. 1901
                  – 12. 12. 1901?] in Westerwede}\toendnotes[C]{\smallbreak}
\Standort{DLA, A:Schnitzler, HS.2022.81.}
\physDesc{Brief, 2 Blätter, 4 Seiten, 1325 Zeichen
\newline{}Handschrift: schwarze Tinte, deutsche Kurrent}
\buchAbdrucke{\weitereDrucke{1) \emph{Rainer Maria Rilke und Arthur Schnitzler. Ihr Briefwechsel.}. Mit Anmerkungen herausgegeben von Heinrich Schnitzler In: \emph{Wort und Wahrheit}, Jg. 13, Nr. 4, April 1958, S. 282–298, hier S. 289.} \weitereDrucke{2) Arthur Schnitzler: \emph{Briefe 1875–1912}. Herausgegeben von Therese Nickl und Heinrich Schnitzler. Frankfurt am Main: \emph{S. Fischer} 1981, S. 439.} }\toendnotes[C]{\smallbreak}
\pstart
           \noindent{}{\pb}mein lieber Herr Rilke, es iſt kaum denkbar,{ }ſchöner über mein
               kleines Stück\pwindex{Schnitzler, Arthur 15.\,5.\,1862 Wien – 21.\,10.\,1931 ebd.@\textsc{Schnitzler, Arthur} (15.\,5.\,1862 Wien – 21.\,10.\,1931 ebd.), \emph{Schriftsteller, Mediziner}!Lebendige Stunden@\strich\emph{Lebendige Stunden}|pwv} auszuſprechen,
               als Sie in Ihrem letzten \label{K_L04344-1v}\edtext{Briefe}{\lemma{\textnormal{\emph{Briefe}}}\Cendnote{\textnormal{XXXX Auszeichnungsfehler: Dokument L04540 nicht gefunden.}}}\label{K_L04344-1} gethan, und ich danke Ihnen
               herzlich für Ihr Verſtehen und für Ihre Freundschaft.\pend
           
\pstart
           Mit der Zuſendung Ihres Buches\pwindex{Rilke, Rainer Maria 4.\,12.\,1875 Prag – 29.\,12.\,1926 Montreux@\textsc{Rilke, Rainer Maria} (4.\,12.\,1875 Prag – 29.\,12.\,1926 Montreux)!Letzten@\strich\emph{Die Letzten}|pwv} haben Sie mir Freude gemacht. Die erſte \label{K_L04344-2v}\edtext{Skizze\pwindex{Rilke, Rainer Maria 4.\,12.\,1875 Prag – 29.\,12.\,1926 Montreux@\textsc{Rilke, Rainer Maria} (4.\,12.\,1875 Prag – 29.\,12.\,1926 Montreux)!Im Gespräch@\strich\emph{Im Gespräch}|pw} kannte ich{ }ſchon aus den Flugblättern\pwindex{Frühling. Moderne Flugblätter. 4. Heft@\emph{Frühling. Moderne Flugblätter. 4. Heft}|pwv}}{\lemma{\textnormal{\emph{Skizze … Flugblättern}}}\Cendnote{\textnormal{Vgl. XXXX Auszeichnungsfehler: Dokument L04343 nicht gefunden.}}}\label{K_L04344-2};{ }ſie{ }ſcheint
               mir {\pb}geiſtreicher als lebendig. Die Novellette »Der Liebende\pwindex{Rilke, Rainer Maria 4.\,12.\,1875 Prag – 29.\,12.\,1926 Montreux@\textsc{Rilke, Rainer Maria} (4.\,12.\,1875 Prag – 29.\,12.\,1926 Montreux)!Liebende@\strich\emph{Der Liebende}|pw}«{ }ſähe ich am liebſten als Einakter.
               Ich glaube es bräuchte nicht viel, um ein{ }ſehr feines Drama draus zu machen.
               Vielleicht nur einige Striche. – Die »Letzten\pwindex{Rilke, Rainer Maria 4.\,12.\,1875 Prag – 29.\,12.\,1926 Montreux@\textsc{Rilke, Rainer Maria} (4.\,12.\,1875 Prag – 29.\,12.\,1926 Montreux)!Letzten [Erzählung]@\strich\emph{Die Letzten [Erzählung]}|pw}«
               haben mir auch direkt einen Fehler, um den Sie \introOben{}von\introOben{} manchen
               beneidet werden könnten. Es iſt zu viel drin – ein letzte{[}s{]}
               Kapitel, in das Sie den Inhalt {\pb}alles des vorhergegangnen,
               verſchwiegnen hineindrängen wollten. So{ }ſind{ }ſie nicht durchaus zur
                  h\textcolor{gray}{öch}ſten Klarheit gekommen. Allerdings war das auch nicht Ihre
               Abſicht.\pend
           
\pstart
           Wie gern möcht ich noch mehr und über allerlei andres mit Ihnen reden – mit Ihnen:
               das heißt: auch gleich Ihre Gegenrede hören. Kommen Sie nicht wieder einmal nach Wien\oindex{Wien@\textbf{Wien}|pw}? {\pb}Oder wenigſtens
               nach Berlin\oindex{Berlin@\textbf{Berlin}|pw} (wo ich Anfang Jänner{ }ſein dürfte?) Auf Ihr \label{K_L04344-3v}\edtext{Stück\pwindex{Rilke, Rainer Maria 4.\,12.\,1875 Prag – 29.\,12.\,1926 Montreux@\textsc{Rilke, Rainer Maria} (4.\,12.\,1875 Prag – 29.\,12.\,1926 Montreux)!tägliche Leben. Drama in zwei Akten@\strich\emph{Das tägliche Leben. Drama in zwei Akten}|pwv} bin ich{ }ſchon durch Berger\pwindex{Berger, Alfred von 30.\,4.\,1853 Wien – 24.\,8.\,1912 ebd.@\textsc{Berger, Alfred von} (30.\,4.\,1853 Wien – 24.\,8.\,1912 ebd.)|pw}}{\lemma{\textnormal{\emph{Stück … Berger}}}\Cendnote{\textnormal{Vermutlich bei ihrem Gespräch am A. S.: \emph{Tagebuch}, 30. 10. 1901.}}}\label{K_L04344-3} geſpannt worden, der
               mir neulich davon geſprochen hat.\pend
           
\pstart
           Ich grüße Sie vielmals und herzlich.{\\[\baselineskip]}Auch Ihrer verehrten Frau\pwindex{Westhoff, Clara 21.\,11.\,1878 Bremen – 9.\,3.\,1954 Fischerhude@\textsc{Westhoff, Clara} (21.\,11.\,1878 Bremen – 9.\,3.\,1954 Fischerhude)|pwv} danke ich für die freundliche
               Theilnahme an meinen Sachen.{\\[\baselineskip]}Ihr \spacefill\mbox{Arth Schnitzler}\pend
           \leftskip=0em{}
\pstart
           Wien\oindex{Wien@\textbf{Wien}|pw}, 7. 12. 901.\pend
           \selectlanguage{ngerman}\endnumbering\briefempfaengerindex{Rilke, Rainer Maria@\textsc{Rilke, Rainer Maria}!zzzSchnitzler, Arthur@\emph{von Arthur Schnitzler}!1901-12-071@{7. 12. 1901}|)be}\mylabel{L04344h}
\begin{anhang}
\end{anhang}\newcommand{\dateiname}{L04344}\newcommand{\titel}{Arthur Schnitzler an Rainer Maria Rilke, 7. 12. 1901}\newcommand{\editorInnen}{Herausgegeben von Jahnke, SelmaMüller, Martin Anton}\input{../tex-inputs/latex-leseansicht-abspann}
